% 【基础正文】所属：第1章（由 chapters/revised/01-knowledge.tex \input 引入）
\minitopic{从“知道”开始，但不止于知识}传统知识论常从命题知识开始。所谓\term{命题知识}{propositional knowledge}，形式是“主体 $S$ 知道命题 $p$”。它不同于熟悉某人、某地的\term{亲知}{knowledge by acquaintance}，也不同于会游泳、会修车的\term{能力之知}{knowledge-how}。这三类用法之间是否能够还原，仍有争议；本书以命题知识为主，但不会假定它穷尽所有认识成就。 知识首先具有\term{事实性}{factivity}：若 $S$ 知道 $p$，则 $p$ 为真。我们可以相信一个假命题，却不能知道一个假命题。事实性不表示人永不犯错，只表示当某次判断确实算作知识时，它没有错。第二，知识通常要求主体持有某种认知态度。把正确答案印在纸上而从未读过它，不会使人知道答案。第三，知识似乎要求信念与真之间不是偶然重合。开篇案例中的乙说中了，却缺乏把真相与判断连接起来的认识根据。

由此出现四组不同的评价：
\begin{enumerate}
  \item \textbf{真值评价}：命题是否为真？
  \item \textbf{认识地位评价}：主体是否知道、是否得到确证？
  \item \textbf{理性评价}：就主体拥有的信息而言，相信它是否合理？
  \item \textbf{责任评价}：主体是否尽到了查证、倾听和修正的责任？
\end{enumerate}
这些评价可能分离。这里把“确证”与“认识理性”暂时分栏，是为了训练诊断，不预设二者在理论上必然不同；有些理论会把它们部分甚至完全重合，另一些理论则用“确证”指更窄的知识相关地位。一个谨慎的人可以根据强证据形成后来被证明为假的信念；一个粗心的人也可能猜中。认识论的任务之一，就是说明这些评价的关系，而不是把它们压缩成“对”与“错”。

\minitopic{问题、对象与层次}认识论至少研究四类问题。第一是\textbf{分析问题}：知识、确证、证据、理解分别是什么？第二是\textbf{来源问题}：知觉、记忆、推理、内省和证言如何产生知识？第三是\textbf{范围问题}：我们能否知道外部世界、他人心灵、道德事实和未来事件？第四是\textbf{规范与社会问题}：应当如何形成信念，制度和权力又如何塑造可知之物？ 把这些问题分开很重要。“人是否有自由意志”主要是形而上学问题；“我们如何知道人有自由意志”才是认识论问题。“某新闻为假”是事实判断；“在当时信息条件下转发它是否不负责任”则是规范判断。问题层次不清，常使争论双方各自回答不同问题。

\minitopic{概念分析与反例}一种经典方法是提出知识的必要条件和充分条件。若 $A$ 是 $B$ 的必要条件，则没有 $A$ 就没有 $B$；若 $A$ 是 $B$ 的充分条件，则有 $A$ 就足以有 $B$。定义“知识是得到确证的真信念”，既声称真、信念和确证各自必要，也声称三者合取充分。

\begin{argumentbox}
概念分析的基本检验：先提出条件集合 $C$；再寻找满足 $C$ 却不是目标概念的案例，以反驳充分性；或寻找属于目标概念却不满足 $C$ 的案例，以反驳必要性。反例不是故事比赛。有效反例必须说明判断稳定、相关细节受控，并指出究竟攻击哪个条件。
\end{argumentbox}

思想实验依赖直觉，但直觉不是神秘证据。我们应检查措辞效应、文化差异、背景假设和概念的实际用途。即使个别直觉分歧，思想实验仍可揭示理论承诺：如果一个理论把明显的猜中算作知识，我们至少知道它对认识运气非常宽容。 除概念分析外，认识论还使用心理学实验、概率模型、语言分析、思想史和社会制度研究。方法的选择应服从问题。有关记忆可靠性的经验事实不能仅靠扶手椅思辨决定；有关“知道”一词语义的结论，也不能直接从脑成像推出。

\begin{comparison}
《论语》“知之为知之，不知为不知”常被理解为一种认识谦逊规范，而不是知识的充分必要条件\parencite[2.17]{analects2003}。若直接把“知”翻译成当代分析哲学的命题知识，就会忽略其修身语境；若因此断言它毫无认识论意义，又会错过它对自我评估和言说责任的关注。比较的第一原则，是先说明原概念在自身实践中的功能，再讨论结构上的可比性。
\end{comparison}

\minitopic{怎样读一篇认识论论文}第一遍阅读只回答三个问题：作者要反对什么看法，作者自己的结论是什么，核心例子是什么。不要在第一遍就逐句查词，否则容易见树不见林。第二遍把论证写成编号前提，补出作者没有明说却必需的桥梁原则。例如，从“该信念由可靠过程产生”到“该信念得到确证”，需要一个把可靠性与确证连接起来的规范前提。第三遍才检验每个前提，并寻找作者最可能接受的修正。 摘要和评价应严格分开。摘要阶段应使观点的支持者也愿意承认自己的立场得到公平呈现；评价阶段则指出反例、歧义或代价。“我不同意”不是反驳，“这个案例不现实”也通常不够，因为思想实验的功能可能只是检验条件组合。有效批评要说明不现实细节为什么破坏了目标直觉或推理。 阅读中还要标记三类词。第一类是作者明确定义的技术词，不能随日常语感替换。第二类是论证连接词，如“因此”“除非”“只要”，它们暴露推理结构。第三类是范围词，如“所有”“通常”“在正常条件下”，许多反例正是通过改变范围起作用。

\minitopic{案例工作坊：天气应用的认识链}设想一款应用显示“明日降雨概率80\%”。用户相信明日会下雨。至少可以提出五种不同问题：（1）模型的输出是否校准；（2）用户是否理解80\%并不等于必然；（3）应用是否基于本地区数据；（4）用户的信念是否由读数形成，而不是原先迷信；（5）带伞行动是否需要达到知识门槛。若第二天未下雨，不能仅凭单次结果断言模型不可靠；若应用伪造数据，即使碰巧下雨，用户也可能缺少知识。 这个案例显示，认识论分析不是给“知道”做一次性判决。我们可以评价模型、信息传递、用户理解和行动选择四个环节。把认识链拆开，才能确定失败发生在哪里，也才能提出相称的改进：校准问题需要统计检验，来源问题需要数据审计，理解问题需要更好的风险表达，行动问题需要结合损失权衡。

\minitopic{常见误区}第一，把“有理由”理解为“能够说出一个理由”。婴儿和动物可能以可靠知觉拥有知识，却不能进行语言辩护。第二，把“主观确信”当作确证程度。确信是一种心理强度，确证是规范地位，两者可能反向变化。第三，把认识论问题全都理解成怀疑论问题；实际研究还关心比较理由、解释理解、制度纠错和认识公正。第四，以为承认社会与历史条件就必须放弃真理。知识生产受条件塑造，与命题是否为真是两个层次；社会分析恰可说明哪些条件改善了发现真相的机会。
